\section{Fitting intervals and frequency control}\label{sec:fitting}

Fix the singular origin, integer $Q$, plane $\Pi$, and target scales
$L_i\downarrow0$ supplied by Proposition~\ref{prop:flat-scales}.
The scales $L_i$ retain the measure of density-$Q$ centers, but were not
chosen by any height or frequency criterion. We shall cover them by
intervals on which a fitted center has uniform height doubling. The
construction of the local fits comes from \cite[Section~7]{AE}; the
selection and comparison of successive fitting intervals are needed here
because the origin need not satisfy a typical-point multiplicity condition.

At each start, local minimal graphs fit the support until their height
test first fails. We glue these graphs to one reference disk. The failure
scales determine the smallest radius on which we estimate its height
and frequency. There are then two possibilities. If this radius is zero,
the frequency estimates apply inward indefinitely. Otherwise we restart
at a smaller scale where the varifold is again close to a plane.
On late intervals carrying the targets, the preceding failure bounds the
frequency at the outer end, and the new failure bounds it at the inner
end. These two endpoint bounds are what make height doubling
uniform. The last subsection proves that the resulting intervals contain
every sufficiently late $L_i$; no coverage of arbitrary small scales is
required.

\subsection{Local fits and their gluing}

A \emph{good start} is a sufficiently small dyadic scale $R$ for which
$V_R$ is within a fixed small weak-varifold distance of a multiplicity-$Q$
plane through the origin in $\mathbf B_{2000}$. Use a metric on varifolds
with bounded mass, and choose a plane whose distance realizes, or
approximates within a fixed factor, the best distance. In particular, if
this distance tends to zero, the corresponding varifolds converge to
their chosen planes after taking a subsequence. Throughout a start we
work in units of $R$ and take its chosen plane as horizontal; $\mu$
denotes the weight of $V_R$ until target units are introduced explicitly.

The goodness threshold will be decreased a finite number of times below.
Weak convergence to a plane, support convergence, and monotonicity then
provide a cylinder with base containing $B_{300}$, fixed vertical room,
and a support-free gap at its vertical boundary. From now on $\mu$
is restricted to this working cylinder, as in
Section~\ref{sec:normal-inputs}; every support condition in the fitting
construction refers to this restriction. On its interior we have
the upper mass bound $\mu(\mathbf B_u(X))\le C u^m$ at all radii used
below, arbitrarily small support height and weak mass error, and small
tilt. The scalar estimate of \cite[Lemma~2.5]{AE} also gives small
projected mass defect on smaller cylinders. These statements concern
all support points in the cylinder, not a selected collection of sheets.
Since the density at the origin is $Q$, we require the start to be small
enough that all needed smaller and fixed-factor larger density ratios
at the origin are uniformly as close to $Q$ as prescribed.

Choose a finite center derivative order $J>m/2+6$, larger than the orders
needed for the signed estimate and subsequent compactness arguments.
Apply the fitting lemma with order at least $J+1$; denote its exponent
$k_0$ by $k$. In particular,
\begin{equation}\label{eq:fitting-order}
 k>4+m/2.
\end{equation}
Next choose its height parameter $m_0>0$ sufficiently small. We recall
precisely the fitting information to be used.

\begin{lemma}[Local fitting inputs]\label{lem:local-fits}
For every $x\in B_5$, there are smooth minimal graphs
$S_{x,\ell}=\operatorname{graph}f_{x,\ell}$ over $B_{80\ell}(x)$,
at descending dyadic scales $\ell=2^{-j}$ starting at $j=0$, through
the first scale at which
\[
 \sup\{|X_v-f_{x,\ell}(X_h)|:
          X\in\spt\mu,\ X_h\in B_{64\ell}(x)\}
                                  \le m_0\ell^k
\]
fails. Every constructed fit, including the first failed one, has
support height at most $Cm_0\ell^k$ on the fixed interior enlargements
of this test ball. On $B_{70\ell}(x)$ its normalized tilt and projected
mass defect are bounded by $Cm_0^2\ell^{2k-2}$.
For a fixed-base smooth vector test $\eta$ supported in
$B_{70\ell}(x)$, put $\eta_{x,\ell}(w)=\eta(x+\ell w)$ on
$B_{70}$. Its first height moment is bounded by
$Cm_0^2\ell^{m+2k-1}\|\eta_{x,\ell}\|_{H^q(B_{70})}$,
for some fixed integer $q$. Thus the test norm is taken after rescaling
the base ball by $\ell^{-1}$, with no additional amplitude factor.
Comparable overlapping fits differ in derivative
order $j\le J+1$ by at most $Cm_0^2\ell^{2k-1-j}$ on smaller common
balls. Their fixed-coordinate derivatives are uniformly bounded, and
their slopes are arbitrarily small when $m_0$ is small.

At a first failure $(x,\ell)$, every fixed sufficiently small comparable
base subball with center in $B_{30\ell}(x)$ contains a set of base
volume at least $c\ell^m$ whose normal fibers contain two entries
separated by at least $cm_0\ell^k$. Both entries have small tangent tilt
relative to the fit and lie at height at most $Cm_0\ell^k$ from it.
The constants may depend on the fixed relative subball radius and tilt
threshold, but not on the start or the failed scale.
\end{lemma}

These are \cite[Lemmas~7.1--7.2 and Equations~(7.7)--(7.8)]{AE}.
Only small initial height and mass error are required for these inputs.
Indeed the scale-zero correction uses the horizontal plane and smooth
mass averages identifying $Q$; subsequent corrections use the preceding
fit and its selected integer. Compatibility and failed-fit separation
use these estimates and the failed height test. The multiplicity-defect
condition imposed later in that source is not used in these arguments.
At every new scale the preceding fit already bounds all support heights
on a horizontal enlargement. Consequently the normal cylinders used
in the estimates retain a support-free vertical gap and sufficient
interior room. The initial cylinder supplies this property at scale zero.

Write $\ell_x$ for the first failed scale, and put $\ell_x=0$ if no
failure occurs. We use the failures in two ways. Their tents prescribe
a spatially varying scale for gluing the local fits; the largest failure
whose $16$-fold base ball contains the origin determines the lower end
of the radial interval. Define these two quantities by
\begin{equation}\label{eq:tents}
 t(w)=\sup_{x\in B_5}(\ell_x-|w-x|)_+,
 \qquad
 b=\max\bigl(\{\ell_x:|x|\le16\ell_x\}\cup\{0\}\bigr).
\end{equation}
The maximum exists because the nonzero scales are dyadic. The function
$t$ is $1$-Lipschitz, $t(0)\le b$, and hence $t(w)\le b+|w|$.
By decreasing the goodness threshold, every failure scale can be made
smaller than any prescribed fixed positive number. Otherwise the
separated entries in Lemma~\ref{lem:local-fits}, whose direction is
nearly vertical, contradict convergence of the support to the start
plane on the initial cylinder.

Glue the fits using the construction in the proof of
\cite[Proposition~7.3]{AE}, before its multiplicity-defect argument.
For clarity, regularize $t$ by $t_\epsilon=\max(t,\epsilon)$, cover by
balls of radii a small fixed multiple of $t_\epsilon$, and choose a
partition of unity with derivatives controlled at those scales. At a
cell center use a constructed dyadic fit of scale comparable to
$t_\epsilon$ and larger than its failure scale. The overlap estimate
in Lemma~\ref{lem:local-fits} controls derivatives of the partitioned
differences. Passing to a $C^J$ limit as $\epsilon\downarrow0$ gives a
graph $S$ with small slope and uniformly bounded $C^J$ geometry, for which
\begin{equation}\label{eq:glued-center}
 |\mathcal H_S|+|\nabla^\perp\mathcal H_S|
       \le Cm_0^2t^{2k-4},
 \qquad |z|\le Cm_0t^k\quad\hbox{on the support}.
\end{equation}
Here $t$ is evaluated at the fixed horizontal argument; in the height
bound one may equivalently use the horizontal argument of the nearest
projection. The same estimate holds for vertical graph-height
differences at fixed horizontal argument. On a cell of size comparable
to $t>0$, the glued graph differs from a constructed fit of comparable
scale by at most $Cm_0^2t^{2k-1}$ in height. These statements follow
from the overlap estimate for the finitely many comparable fits in
each partition sum; they require no density statement about $\{t>0\}$.

We make slopes, normal tubes, and geodesic coordinate neighborhoods
small enough that the relevant parametrizations and nearest projections
have Lipschitz constants below $3/2$. The point of $S$ over horizontal
$0$ is at distance at most $Cm_0b^k$ from ambient $0$. Use this point
as radial center on $S$, and denote geodesic distance from it by $\rho$.
Fix $B_*=100$ and a sufficiently small $r_*>0$. We shall work on
\begin{equation}\label{eq:fitting-radius-range}
 B_*b\le r\le4r_*,\qquad r>0.
\end{equation}
All these radii lie strictly inside one geodesic chart. The maximum
failure scale is taken small enough that the interval has the required
interior room.

\subsection{Weighted height on a fitting interval}

Using the integer $p\ge4$ fixed in Section~\ref{sec:normal-inputs}, define
\[
 \phi=(1-\rho^2/r^2)_+^p,
 \qquad \psi=2p(1-\rho^2/r^2)_+^{p-1}.
\]
Use the normal-cylinder notation of the preceding section. In particular,
$B=I-(P_T)_{hh}$, $e=\tr B$, $J_z=I-\mathcal S_z$, and $K_z=J_z^{-1}$.
Put $u=\nabla_S\rho$, $v_z=K_zu$, and
\[
 b_z=B[v_z,v_z],\qquad h=1-b_z/|v_z|^2.
\]
Then $0\le h\le1$. On the central fiber choose any horizontal unit
vector in place of $u$. If this fiber has positive mass, its tangent is
purely vertical almost everywhere on the fiber by rectifiability, so
$h=0$ there almost everywhere. The reason for this factor is the sharp
projection inequality
\[
 |(P_T)_{vh}v_z|^2
   =B[v_z,v_z]-|Bv_z|^2
   \le b_z-\frac{b_z^2}{|v_z|^2}=b_zh.
\]
Weighting squared height by $h$ will therefore pair it directly with
the radial tilt $b_z$ in Cauchy--Schwarz. No positive lower bound for
$h$ is assumed. Define
\begin{equation}\label{eq:interval-quantities}
 \begin{aligned}
 H&=r^{-m}\int\psi|z|^2h\,\dd\mu,
 &\delta&=r^{-1}\sup_{\rho<r,\ X\in\spt\mu}|z|,\\
 A&=r^{2-m}\int\phi\,\dd M,
 &E_*&=r^{2-m}\int\phi\,\dd D,\\
 R_*&=r^{-m}\int\rho^2\psi b_z\,\dd\mu.
 \end{aligned}
\end{equation}
The radius in the supremum refers to the projected point. These
quantities will compare neighboring radii without requiring that the
center be minimal.

\begin{lemma}[Weighted bounds]\label{lem:weighted-bounds}
On \eqref{eq:fitting-radius-range}, $\delta>0$, $\delta\le Cr^{k-1}$,
and normalized mass is bounded. For $J_0=m+p+4$ as fixed in
Section~\ref{sec:normal-inputs},
\begin{align}
 c\delta^{J_0}&\le H/r^2\le C\delta^2,
                                  \label{eq:height-power-bounds}\\
 |\mathcal H_S|+|\nabla^\perp\mathcal H_S|
   &\le C\delta r\quad(\rho<r),\label{eq:point-curvature-bound}\\
 \int\phi\bigl(|\mathcal H_S|^2+
                    |\nabla^\perp\mathcal H_S|^2\bigr)\,\dd\mu
   &\le C\int\phi|z|^2\,\dd\mu.\label{eq:weighted-curvature-bound}
\end{align}
Moreover a nonnegative smooth test $\chi_r$, with reference integral
one and support in $\{\rho<r/4\}$, can be chosen so that its
projected mass average is within $1/8$ of $Q$ and
$r^m\|\chi_r\|_{C^{J_{\mathrm{sc}}}_r}\le C$.
All constants are uniform among good starts.
\end{lemma}

\begin{proof}
If $t(w)>0$ and $\rho(w)<r$, choose a failure $(x,\ell)$ for which
$\ell-|w-x|>t(w)/2$. Since $|w|\le r$, either $|x|\le16\ell$ or
$16\ell<r+\ell$. Thus
\begin{equation}\label{eq:witness-size}
 \ell\le\max(b,r/15)\le r/15.
\end{equation}
Equation~\eqref{eq:glued-center} gives $\delta\le Cr^{k-1}$.
The projected cylinder is contained in an ambient ball of comparable
radius, so monotonicity gives its normalized mass bound.

If $\delta=0$, the support lies on $S$ in a neighborhood of ambient
$0$: its projection lies inside $\{\rho<r\}$ because its distance
from the radial center is $O(m_0b^k)$.
Tangential stationarity then makes the multiplicity, extended by zero
on the disk, constant. It is positive because $0$ belongs to the
support. Normal stationarity and elliptic regularity make $S$ smooth
and minimal there. This contradicts the singularity of $0$.

We next locate the height that pays for a witness. The argument is the
weighted construction of \cite[Proposition~7.4]{AE}, with
\eqref{eq:witness-size} replacing its typical-point estimate.
If the witness is near the kernel boundary, move inward by $\ell$
along the radial geodesic and take a ball of radius a sufficiently
small fixed multiple of $\ell$; otherwise take a nearby such ball.
This ball lies within the kernel support, and its center remains within
$B_{30\ell}(x)$ in the fitting coordinates. On it $\phi$ and $\psi$
dominate, up to fixed constants, their suprema on a small cell containing
the original point, and
$\psi\ge c(\ell/r)^{p-1}$.
Lemma~\ref{lem:local-fits} supplies separated entries above a set of
base volume $c\ell^m$. Two entries separated in a nearly vertical
direction cannot both lie within a sufficiently small multiple of their
separation from a small-slope graph. After transferring to $S$, at
least one entry therefore has height at least $cm_0\ell^k$ and
$h\ge3/4$. Small slopes and the height bounds make changes of horizontal
argument smaller than the available subball margin. The projection
Jacobian bound converts base volume into the same lower bound for mass.

In particular $\delta r\ge cm_0\ell^k$. At a point approaching the
height supremum, the witness also satisfies
$\delta r\le Cm_0\ell^k$, by \eqref{eq:glued-center}. Hence
\[
 H/r^2\ge c\delta^2(\ell/r)^{m+p-1},
 \qquad \ell/r\ge c\delta^{1/k},
\]
where the last inequality uses $r\le1$. This proves the lower bound
in \eqref{eq:height-power-bounds}, for example with
$J_0\ge2+(m+p-1)/k$; its upper bound follows from the mass bound.
The comparison $\delta r\ge cm_0\ell^k$, together with
$k>4$, proves \eqref{eq:point-curvature-bound}.

To prove \eqref{eq:weighted-curvature-bound}, group witnesses by dyadic
scale and use each covering grid cell once at its scale. By
\eqref{eq:glued-center}, the curvature-square integral on a cell of
size $\ell$ is at most
$Cm_0^4\ell^{4k-8+m}\sup\phi$.
Its paying ball contributes at least
$cm_0^2\ell^{2k+m}\sup\phi$ to the height-square integral.
At each scale the paying balls have bounded overlap, since they are
moved by only $O(\ell)$. Summing the factors
$Cm_0^2\ell^{2k-8}$ proves the assertion. The cells may be used as an
overcover, so no measurable choice of witnesses is necessary.

Finally choose a smooth decreasing radial mass test on a small fixed
fraction of $r$. Compare it with the corresponding ambient radial test
at $0$. The displacement of the center is $O(m_0b^k)$, the support
height is $O(r^k)$, and chord and geodesic radii on $S$ differ by
$O(r^3)$. The radius replacement has error $O(r^2)$, so the normalized
averages differ by $O(r)$, including their reference normalizations.
All support encountered by the ambient test lies in the cylinder and
has projected radius below $r/2$. The ambient average is close to $Q$
by the density-ratio requirement at the start. Decreasing the fixed
parameters gives the stated $1/8$ bound.
\end{proof}

\subsection{Frequency estimates and their inward consequences}

The preceding lemma gives positive height and the single-radius
hypotheses for the frequency calculation. Set
\[
 P=A/H,\qquad n_*=E_*/H,\qquad q_*=R_*/H,
 \qquad T=\log(r^2/H),\qquad g=T^{-20}.
\]
Then $T$ is comparable to $|\log\delta|$ and
$T\ge2|\log r|-C$. A dot denotes differentiation with respect to
$\log r$.

The use of signed mass in $P=A/H$ is dictated by horizontal stationarity,
which controls $\dot A$. To see the intended cancellation, temporarily
omit the geometric errors in the estimates below. Signed excess gives
$2P\le n_*$, the sharp flux inequality gives $n_*^2\le q_*$,
and the variation identities give
$\dot A/H=2P-n_*+q_*$ and $\dot H/H\le2n_*$.
For $P\ge0$ these imply
\[
 \dot P\ge2P-n_*+n_*^2-2Pn_*
             =(n_*-1)(n_*-2P).
\]
This expression is nonnegative when $P\ge1/2$.
The next two lemmas quantify the errors and recover the inward bounds
and height comparison needed on a fitting interval.

\begin{lemma}[Frequency on a fitting interval]\label{lem:interval-frequency}
After fixing $r_*$ sufficiently small, on
\eqref{eq:fitting-radius-range} the following inequalities hold:
\begin{align}
 |P|&\le C(n_*+r^2)+\eta q_* ,\label{eq:freq-absolute}\\
 2P&\le n_*+Cg(n_*+1)+Cr^2+\eta q_* ,\label{eq:freq-signed}\\
 n_*&\le\sqrt{q_*}+Cr^2+\eta n_* ,\label{eq:freq-flux}\\
 \left|\dot A/H-(2P-n_*+q_*)\right|
 &\le C\{r\sqrt{q_*}+r^2(1+n_*)\}+\eta(n_*+q_*),
                                      \label{eq:freq-horizontal}\\
 -m\le\dot H/H&\le2n_*+Cr+\eta(n_*+q_*).
                                      \label{eq:freq-height}
\end{align}
Here $\eta=C\delta^\alpha\le C'e^{-cT}$ for fixed positive
$\alpha,c$, and constants are uniform among intervals. The functions
whose derivatives occur are locally absolutely continuous in $\log r$.
\end{lemma}

\begin{proof}
We verify the transfer of the calculation in
\cite[Lemmas~8.1--8.2]{AE}; this also identifies the precise
single-radius inputs. The definitions and positivity of $B$ give
\begin{equation}\label{eq:weighted-height-comparisons}
 \begin{aligned}
 r^{-m}\int\phi|z|^2\,\dd\mu&\le C(H+\delta^2E_*),\\
 r^{-m}\int\psi|z|^2\,\dd\mu
     &\le C\{H+\delta^2(E_*+R_*)\},\\
 r^{-m}\int\phi\,\dd\Lambda&\le C(H+\delta^2E_*).
 \end{aligned}
\end{equation}
For the first bound use $1-h\le Ce$; for the second split at $r/2$
and use $1-h\le Cb_z$ on the outer part. The last bound also uses
\eqref{eq:weighted-curvature-bound}, with physical units equal to the
current start units.

Here are the details of the mass comparison, where preserving the
coefficient one in \eqref{eq:freq-signed} is essential. For the fixed
$0<\beta_1<1/4$ put $\gamma_1=\delta^{\beta_1}$. Away from the
outer collar we use signed excess; on the collar we use a radial
first-variation test. Choose a radial cutoff $\chi$ equal to one at
relative depth $1-\rho/r\le\gamma_1^2$ and zero at depth at least
$2\gamma_1^2$, and put
\[
 \phi_{\rm thin}=\chi\phi,\qquad
 \phi_{\rm bulk}=(1-\chi)\phi.
\]
Cover the support of $\phi_{\rm bulk}$ by cells centered at $x_j$ of
radius $\lambda_1=c_1\gamma_1^2r$, with a smooth partition of unity
$\vartheta_j$ and bounded overlap of fixed enlargements. Choose $c_1$
small enough that these enlargements stay inside $\{\rho<r\}$ and
$\phi$ is comparable there to $\phi_j=\phi(x_j)$. Their mass is
$O(\lambda_1^m)$ because
$\delta r=o(\lambda_1)$. Testing normal stationarity with a cutoff
times $z$ and using \eqref{eq:point-curvature-bound} yields
\[
 \lambda_1^{-m}(D+\Lambda)(B_{c\lambda_1})
                    \le C\delta^{2-4\beta_1}.
\]
This is the local estimate in \cite[Equation~(8.6)]{AE}.
The scalar estimate selects an integer on each cell; the small
exceptional projection sets cannot fill the prescribed overlaps, so
the integers agree throughout the connected bulk; call their value $q$.
Put $\epsilon=C\delta^{2-4\beta_1}$.
Partition the smooth mass test from Lemma~\ref{lem:weighted-bounds}
over the cells. Each piece has the scaled derivative norm in
Equation~\eqref{eq:scalar-distribution-before-count} bounded by
$Cr^{-m}$. That homogeneous estimate bounds its error relative to $q$
by $C\epsilon\lambda_1^m r^{-m}$, without dividing by the piece's
possibly small integral. Bounded overlap gives
$\sum\lambda_1^m\le Cr^m$, so the total error is $O(\epsilon)$.
Since the test average is within $1/8$ of $Q$, integrality gives $q=Q$.

The scalar estimate now bounds $|M|(B_{4\lambda_1}(x_j))$ by
$C(D+\Lambda)(B_{8\lambda_1}(x_j))$. Summing with comparable
weights gives the absolute bulk estimate
\begin{equation}\label{eq:bulk-absolute}
 r^{2-m}\int\phi_{\rm bulk}\,\dd|M|
 \le Cr^{2-m}\int\phi\,\dd(D+\Lambda)
 \le C(E_*+r^2H).
\end{equation}
In the last inequality a small multiple of $E_*$ is absorbed using
\eqref{eq:weighted-height-comparisons}.

For signed excess, use fresh charts on the cells and set
\[
 d_j=\max\{C\lambda_1^{-m}\tau(B_{4\lambda_1}(x_j)),
                          \delta^{2J_0+4}\},
 \qquad T_j=-\log d_j,\qquad \tau=|M|+D+\Lambda.
\]
The local bounds and the floor give
$c|\log\delta|\le T_j\le(2J_0+4)|\log\delta|$.
The physical smallness hypothesis follows from
\[
 \lambda_1+\sup|z|\le C\delta^{2\beta_1},
 \qquad \kappa(2J_0+4)<2\beta_1.
\]
These are precisely the fixed parameter choices in
Equation~\eqref{eq:parameter-order}; the required derivative orders
were chosen after $J_0$ and $\kappa$, and before the fitting exponent.
The partitioned tests $\vartheta_j\phi_{\rm bulk}/(C\phi_j)$ have
size and Lipschitz constant at most one in cell units. Applying signed
excess to these tests and restoring the measure factors gives
\begin{align*}
 2r^{2-m}\int\phi_{\rm bulk}\,\dd M
 &\le r^{2-m}\int\phi_{\rm bulk}\,\dd D
     +C_Qr^{2-m}\int\phi_{\rm bulk}\,\dd\Lambda\\
 &\quad+Cr^{2-m}\sum_j\phi_j\lambda_1^m d_j e^{-F_Q(T_j)}.
\end{align*}
The weighted budgets satisfy
\[
 \sum_j\phi_j\lambda_1^m d_j
 \le C\int\phi\,\dd(D+\Lambda)+Cr^m\delta^{2J_0+4},
\]
by scalar control on the enlarged cells and bounded overlap. The
exponential factor is smaller than any fixed negative power of $T$,
and the floor costs
$r^2\delta^{2J_0+4}\le C\delta^{J_0+4}H$ by
\eqref{eq:height-power-bounds}. The signed-excess remainder is therefore at most
$Cg(E_*+H)$. The separate $C_Q\Lambda$ term is bounded by
$Cr^2H+Cr^2\delta^2E_*$; since $r^2\delta^2=o(g)$, it contributes
only $Cr^2H+CgE_*$. This proves the bulk signed comparison with
leading tilt coefficient one.

For the thin part, let $a$ be a multiplier with
$\sup(|a|+r|\nabla_Sa|)\le1$. The radial primitive constructed in
\cite[Lemma~8.1, Equation~(8.11)]{AE}, cut off at relative depths
between $\gamma_1$ and $2\gamma_1$, gives a compactly supported
tangent field $X$ with
\[
 \diver_SX=a\phi_{\rm thin}+q_c,\qquad
 |X|/r+|\nabla^SX-a\phi_{\rm thin}u\otimes u|
                                    \le C\gamma_1^2\phi.
\]
Here $q_c$ is supported on that thicker cutoff collar and
$|q_c|\le C\gamma_1^{p+1}\phi$. Its mass contribution is controlled
by the absolute bulk estimate \eqref{eq:bulk-absolute}. In horizontal
stationarity, the principal derivative of $X$ pairs with radial tilt
$B[u,u]$; the remaining derivative has the small factor $\gamma_1^2$.
The resulting bound is
\[
 \begin{aligned}
 r^{2-m}\left|\int a\phi_{\rm thin}\,\dd M\right|
 &\le Cr^{2-m}\int\phi_{\rm thin}B[u,u]\,\dd\mu
                  +C\delta^\alpha E_*+Cr^2H
 \\
 &\le C\delta^{\alpha'}(E_*+R_*)+Cr^2H.
 \end{aligned}
\]
For the last inequality use, on the thin collar,
$\phi_{\rm thin}\le C\gamma_1^2(\rho/r)^2\psi$ and
$B[u,u]\le2b_z+C|z|^2e$.
The errors use only horizontal stationarity, fixed geodesic
geometry, and \eqref{eq:weighted-height-comparisons}.
Combining the absolute bulk and thin estimates proves
\eqref{eq:freq-absolute}, including its version with any bounded
scaled $C^1$ multiplier. For the signed estimate take $a=1$ and
combine with the signed bulk bound. The small tilt error is absorbed
in $CgE_*$, since every positive power of $\delta$ is $o(g)$; the
radial error remains $\eta R_*$. This proves \eqref{eq:freq-signed}.

Normal stationarity gives the flux estimate
\[
 \left|E_*-r^{-m}\int\rho\psi\,
                    z\cdot(P_T)_{vh}v_z\,\dd\mu\right|
                       \le Cr^2H+\eta E_*.
\]
Projection algebra gives the exact inequality
$|(P_T)_{vh}v_z|^2\le b_zh$; Cauchy--Schwarz proves
\eqref{eq:freq-flux}. Horizontal stationarity with the lift of
$\rho u$ gives \eqref{eq:freq-horizontal}, as in
\cite[Lemma~8.2]{AE}. Its geodesic volume-distortion multiplier has
bounded scaled $C^1$ norm, so the preceding absolute estimate applies
with a uniform constant.

One modification is needed for the height identity because our center
need not pass through ambient $0$. Off the central fiber put
\[
 Y=\frac{\rho K_zu}{|K_zu|^2}.
\]
Then $DY=P_h+O(\rho+|z|)$, where $P_h$ is horizontal projection
at the nearest point of $S$. To see this, the pulled-back radial field
$\rho u$ has derivative $P_h+O(\rho+|z|)$. Its correction in $Y$
is $\rho$ times a function of the tube coordinates and $u$ that is
$O(|z|)$, including its angular derivatives. Angular differentiation
therefore costs $\rho\cdot O(|z|)/\rho=O(|z|)$, and the remaining
derivative costs $O(\rho)$.
Testing with $\psi|z|^2Y$ gives the following identity, whose integrand
is extended by zero on the central fiber:
\begin{equation}\label{eq:interval-height-identity}
 \dot H=r^{-m}\int\psi\left\{
 |z|^2(\tr(P_TDY)-mh)
 +\frac{2\rho z\cdot(P_T)_{vh}K_zu}{|K_zu|^2}
 \right\}\,\dd\mu .
\end{equation}
This is \cite[Equation~(8.18)]{AE}. First omit the central fiber
using a radial cutoff. The transition term tends to zero because
$Y=O(\rho)$ and the mass of the open shrinking transition annulus
tends to zero. On the fiber $h=0$ almost everywhere, so the left-hand
side is unchanged; the flux term has a factor $\rho$. The error in
$DY$ is $O(r)$ on the support, and
\eqref{eq:weighted-height-comparisons} and the normal flux estimate
give the upper bound in \eqref{eq:freq-height}. The lower bound follows
by direct differentiation of the nonnegative weight $r^{-m}\psi$.
All parameter choices have been fixed independently of the start.
\end{proof}

We now extract the ODE consequences before selecting successive starts.
The proof adapts \cite[Proposition~8.3]{AE} to finite intervals with
controlled endpoint values.

\begin{lemma}[Inward control]\label{lem:inward-control}
There is a fixed large $K$ such that, writing $v_0=\dot A/H$,
\begin{equation}\label{eq:large-frequency}
 |P|\ge K\quad\Longrightarrow\quad
 q_*/|P|\ge c\min(|P|,\eta^{-1}),\qquad v_0\ge cq_*.
\end{equation}
When $K\le P\le T^4$,
\begin{equation}\label{eq:moderate-frequency}
 \dot P\ge-C(rP+gP^2).
\end{equation}
A lower bound for $P$ at the smaller-radius endpoint of a finite
interval gives a lower bound throughout that interval. On an interval
extending to zero, $P$ is bounded below for all sufficiently small radii.
If at a sufficiently small starting radius $r_0$ one has $P_0\le T_0^2$,
then inward throughout the interval
\begin{equation}\label{eq:inward-upper}
 P\le \exp(1)\max(K,P_0).
\end{equation}
Such a starting radius exists on every interval extending to zero.
Finally, wherever $|P|\le C_0$ and $\eta$ is sufficiently small
depending on $C_0$, every fixed bounded-ratio subinterval satisfies
\begin{equation}\label{eq:height-doubling}
 H(r')\asymp H(r''),\qquad
 \int q_*\,\dd\log r\le C.
\end{equation}
The constants in the last conclusion depend on $C_0$ and the radius
ratio, not on the position of the interval.
\end{lemma}

\begin{proof}
Equations~\eqref{eq:freq-absolute} and \eqref{eq:freq-flux} imply
$|P|\le C(\sqrt{q_*}+r^2)+C\eta q_*$. For $|P|$ large, either
the square-root term or the last term controls a fixed fraction of
$|P|$, proving the first assertion of \eqref{eq:large-frequency}.
Equation~\eqref{eq:freq-horizontal} then gives its second assertion,
after increasing $K$ and decreasing $r_*$.

For positive $P$, quotient differentiation and the upper height bound
give the principal expression
$2P-n_*+q_*-2Pn_*$, with the errors displayed in
Lemma~\ref{lem:interval-frequency}. If $q_*$ is larger than a
sufficiently large fixed multiple of $P^2$, it absorbs all negative
terms, since $n_*\le C\sqrt{q_*}+Cr^2$ and $\eta T^4=o(1)$.
Otherwise $n_*\le CP$ and
\[
 q_*\ge n_*^2-Cr^2P-C\eta P^2,
 \qquad n_*\ge2P-C(gP+r^2+\eta P^2).
\]
For large $K$ the principal terms factor as
$ (n_*-1)(n_*-2P)$, bounded below by
$-CP(gP+r^2+\eta P^2)$. Since $\eta T^4\le g$ at small radii,
the term $\eta P^3$ is absorbed by $gP^2$. This proves
\eqref{eq:moderate-frequency}.

If $P\le-K$, the lower height bound gives
$\dot P=v_0-P\dot H/H\ge cq_*-m|P|>0$.
Thus $P$ decreases in inward time $-\log r$ whenever it is below
$-K$. This proves the finite-interval lower bound from the inner
endpoint. On an infinite interval, such an excursion would also make
$A<0$ decrease inward and remain below a fixed negative value. That is
impossible because $A\to0$: its normalized mass integral is bounded
and it has a factor $r^2$.

To prove the upper bound, work on a component where $P\ge K$ before
any exit through $P=T^4$. In inward time, \eqref{eq:moderate-frequency}
gives
\[
 (\log P)'\le C(r+T^{-16}).
\]
The integral of this right-hand side from a sufficiently small radius
to zero is less than one, since $T\ge2|\log r|-C$. Moreover
$Z=T-\log P=\log(r^2/A)$ satisfies
$Z'=v_0/P-2>0$ by \eqref{eq:large-frequency}.
If the component starts at $r_0$ with $P_0\le T_0^2$, then
$P\le\exp(1)P_0$ and $T\ge T_0-\log P_0+\log K\ge T_0/2$.
For large $T_0$ these inequalities exclude $P=T^4$.
Components starting later at $P=K$ satisfy the same argument, giving
\eqref{eq:inward-upper}.

On an infinite interval, suppose no sufficiently small radius satisfied
$P\le T^2$. Then eventually $P>T^2$, and
\eqref{eq:large-frequency} implies $Z'\ge cT^2-2$.
The lower bound for $T$ makes $Z$ positive eventually. Since $Z\le T$
in this regime, one then has $Z'\ge c'Z^2$, which forces blow-up in
finite inward time. The continuous positive quantities $A,H$ are
defined at every finite time in this regime, a contradiction.

Finally, for $|P|\le C_0$, use the upper height derivative bound if
$P\ge0$ and the lower bound if $P<0$. Absorbing $\eta q_*$ gives
\[
 \dot P\ge cq_*-C_{C_0}.
\]
Integration bounds $\int q_*\,\dd\log r$ on every bounded-ratio
interval. Equation~\eqref{eq:freq-flux} and Cauchy--Schwarz bound
$\int n_*\,\dd\log r$, and \eqref{eq:freq-height} gives both
height comparisons in \eqref{eq:height-doubling}.
\end{proof}

\subsection{Covering the prescribed target scales}

We have obtained height comparison whenever the frequency can be
bounded above and below on a radial interval. For a finite interval,
Lemma~\ref{lem:inward-control} uses an upper bound at its outer endpoint
and a lower bound at its inner endpoint. The restart construction supplies
them from different failures: the preceding start's failure supplies the
outer bound, and the current start's failure supplies the inner bound.
If no failure occurs near the origin, the interval extends to zero and
the infinite-interval conclusions of that lemma apply instead.

We first define the starts and prove that their intervals contain the
prescribed targets. Fix $N_*=1000$, and choose a dyadic factor
\begin{equation}\label{eq:restart-factor}
 F>10N_*B_*/r_*.
\end{equation}
Decrease the maximum failure scale so that
$Fb<1/4$ and $8N_*B_*b<r_*$ for every good start.
Its target interval is
\begin{equation}\label{eq:target-interval}
 B_*bR\le L\le(r_*/N_*)R,\qquad L>0.
\end{equation}
Starting at one good dyadic $R$, if $b>0$ choose the next start to be
the largest good dyadic scale at most $FbR$; if $b=0$, stop. Such a
next start always exists: every dyadic scale comparable to a sufficiently
late $L_i$ by fixed factors is good, because $V_{L_i}\to Q|\Pi|$
locally.

\begin{proposition}[Uniform comparison at the targets]\label{prop:target-frequency}
All sufficiently late $L_i$ belong to intervals
\eqref{eq:target-interval}. After passing to a subsequence, either one
interval extends to zero and contains all these targets, or the indices
of their covering intervals tend to infinity. In the latter case,
on each sufficiently late covering interval,
\begin{equation}\label{eq:uniform-frequency-range}
 |P(r)|\le C\qquad(B_*b\le r\le r_*),
\end{equation}
and \eqref{eq:height-doubling} holds there with uniform constants.
In the former case the same conclusions hold on the sufficiently
small-radius part of the fixed interval.
\end{proposition}

\begin{proof}
\emph{Coverage of the targets.}
The top of the interval belonging to $R$ is $(r_*/N_*)R$.
Successive tops decrease to zero unless the construction stops, since
the next start is at most $FbR<R/4$. Suppose a target $L_i$ lies below
one top and above the next, but misses the first interval. It must then
lie below its lower endpoint, so $L_i<B_*bR$ and
$FbR>10N_*L_i/r_*$. A dyadic scale between
$2N_*L_i/r_*$ and $4N_*L_i/r_*$ is therefore eligible for the next
start. For all sufficiently late $i$ it is good, by convergence at
$L_i$ to a plane. Maximality of the next start would make its top
larger than $L_i$, a contradiction. Hence all sufficiently late targets
are covered.

A fixed interval containing targets tending to zero necessarily has
$b=0$. On that interval Lemma~\ref{lem:inward-control} first supplies
a lower bound and a sufficiently small radius with $P\le T^2$, then
an upper bound and height doubling inward. This proves the assertion
in the fixed-interval case. After taking a subsequence, we may therefore
restrict to targets whose covering interval indices tend to infinity.
There are then infinitely many starts, all with positive failure
parameter $b$.

\emph{Geometry of the late starts.}
We first show that the $b$ values of these covering intervals tend to
zero. Otherwise their start scales would be comparable to their target
scales, since
$B_*b\le L_i/R\le r_*/N_*$. The start varifolds would then converge
to planes, contradicting failed-fit separation at a scale bounded
away from zero.

In fact the start varifolds converge to multiplicity-$Q$ planes even
without comparability to the targets. Any subsequential limit is a
stationary tangent cone at $0$, of density $Q$. Pass also to limits of
the start coordinates and their uniformly controlled center graphs.
Because $t(w)\le b+|w|$ and $b\to0$, the limiting cone lies within
$C|w|^k$ of a $C^1$ graph through $0$. Every conical ray is therefore
in the tangent plane of that graph at $0$. Constancy gives a
multiplicity-$Q$ plane. Thus, if a minus sign denotes the preceding
interval, eventually
\begin{equation}\label{eq:exact-restart}
 R=Fb_-R_-.
\end{equation}
Indeed the right-hand side is dyadic. If it exceeded the chosen $R$,
then $2R$ would still be eligible; planar convergence of $V_R$ on
every fixed ball makes $2R$ good, contradicting maximality. It follows
also that $b_-\to0$. If not, the preceding starts would be comparable
to the current planar starts, again contradicting separation.

\emph{The outer endpoint: the preceding failure.}
We now work entirely in the current start's units. The preceding failure
will give an upper support-height bound and a matching lower height
integral on fixed room. Together these bounds will control $P(r_*)$.
Define the preceding failure height in current units by
\[
 h_0=\frac{m_0b_-^k}{Fb_-}.
\]
In every required fixed bounded region of the current start, all
support points are within $Ch_0$ of the preceding center, expressed
in current units: in preceding units the region has radius $O(Fb_-)$,
and \eqref{eq:glued-center} applies with $t\le b_-+|w|$.
The preceding and current start planes make a small angle. To verify
this, take a subsequence: the preceding centers in current units have
small slope and pass within $Ch_0$ of $0$, and the current planar
support limit lies in their limiting graphs. The current chosen
planes converge to this support plane. This proves the asserted angle
bound, allowing a fixed enlargement of the small slope threshold.

The current center consequently satisfies
\begin{equation}\label{eq:restart-height-upper}
 |z|\le Ch_0
\end{equation}
on the region containing $\{\rho<4r_*\}$. Indeed a current witness
of scale $\ell$ supplies two entries separated by $cm_0\ell^k$ in
a direction nearly vertical also for the preceding plane. Both lie
within $Ch_0$ of its small-slope center, so $cm_0\ell^k\le Ch_0$.
Equation~\eqref{eq:glued-center} now proves
\eqref{eq:restart-height-upper}. By Lemma~\ref{lem:weighted-bounds}
at radius $4r_*$, the mean curvature and its first normal derivative
are at most $Ch_0$ there as well. Every comparison is internal to both
cylinders: the current cylinder has fixed room, whereas its required
region in preceding units shrinks with $Fb_-\to0$.

A preceding failure attaining $b_-$ supplies the complementary bound
\begin{equation}\label{eq:restart-height-lower}
 H(r_*)\ge ch_0^2.
\end{equation}
Its center is within $16b_-$ of horizontal $0$, so the failed-fit
lemma may be used on a subball centered near horizontal $0$.
The entries lie within ambient radius $Cb_-$ in preceding units:
their heights over the fit are $O(m_0b_-^k)$, and the glued center
has the same bound at $0$. In current units they are well inside
$\{\rho<r_*\}$ by \eqref{eq:restart-factor}; their mass is bounded
below by a constant times $F^{-m}$. Their tilt remains small relative
to the current center. At least one entry of each separated pair has
height at least $ch_0$ over that center, proving
\eqref{eq:restart-height-lower}.

The cutoff normal test with $z$ on fixed room around $r_*$, using
\eqref{eq:restart-height-upper} and its curvature bound, gives
$E_*(r_*)+R_*(r_*)\le Ch_0^2$. Hence
\eqref{eq:freq-absolute} and \eqref{eq:restart-height-lower} bound
$|P(r_*)|$ uniformly, while $T(r_*)\to\infty$.
In particular $P(r_*)\le T(r_*)^2$ for all sufficiently late covering
intervals, as required for the inward upper bound.

\emph{The inner endpoint: the current failure.}
For $b>0$ we use the same height-versus-tilt comparison at radius
$B_*b$, now supplied by the current failure. On a fixed enlargement,
$|z|\le Cm_0b^k$ follows from $t\le b+|w|$; the curvature bounds of
Lemma~\ref{lem:weighted-bounds} apply there. The cutoff tilt estimate
on room of scale $b$ gives
\[
 E_*(B_*b)+R_*(B_*b)\le Cm_0^2b^{2k}.
\]
The inverse-square scale in that estimate is canceled by the factors
in \eqref{eq:interval-quantities}. A failure attaining $b$, tested on
a subball of radius $10^{-2}b$ near horizontal $0$, supplies low-tilt
entries well inside the kernel and proves
$H(B_*b)\ge cm_0^2b^{2k}$. Thus $|P(B_*b)|$ is bounded as well.

The outer bound controls $P$ from above as the radius decreases.
The inner bound controls it from below: once $P<-K$, it decreases
inward, so a value below the inner-endpoint bound could not recover
before that endpoint. Lemma~\ref{lem:inward-control} therefore gives
\eqref{eq:uniform-frequency-range}. Its final smallness requirement
on $\eta$ is uniform on late intervals: use $\delta\le Cr$ on
the sufficiently small radii and
\eqref{eq:restart-height-upper} on the remaining radii.
This proves the uniform height comparison.
\end{proof}
